\begin{table}[H]
\centering
\begin{tabular}{lrrrr}
\toprule
& \multicolumn{4}{c}{Auxiliary scale $\eta$} \\
\cmidrule(lr){2-5}
$\lambda$ & $0.4$ & $0.6$ & $0.85$ & $0.95$ \\
\midrule
$0.05$ & $-4.369$ & $-3.686$ & $-3.228$ & $-3.198$ \\
$0.1$ & $-4.287$ & $-3.510$ & $-3.154$ & $\mathbf{-3.133}$ \\
$0.2$ & $-4.381$ & $-3.811$ & $-3.603$ & $-3.614$ \\
$0.4$ & $-5.312$ & $-4.979$ & $-4.855$ & $-4.874$ \\
$0.8$ & $-6.003$ & $-9.129$ & $-9.836$ & $-9.877$ \\
\bottomrule
\end{tabular}
\caption{Two-state calibration grid at $T=5$ with exact population flow: final validation objective, averaged over five seeds, for every pair of perturbation scales. The optimum is interior in $\lambda$, as the bias-variance trade-off of the main perturbation predicts. In $\eta$ the objective increases over the whole tested range whenever $\lambda\leq0.1$ and peaks at $\eta=0.85$ for $\lambda\in\{0.2,0.4\}$, so the auxiliary estimator wants a radius near the top of the grid: at the auxiliary sample sizes used, the $(n\eta^2)^{-1}$ variance term of Proposition~\ref{prop:auxiliary-sensitivity-error} dominates its $O(\eta)$ bias. The ordering reverses only at $\lambda=0.8$, where the run is already dominated by perturbation bias. This grid fixes $\eta=0.85$ for the remaining benchmarks.}
\label{tab:twostate-grid}
\end{table}
